\documentclass[conference]{IEEEtran}

% Required Packages
\usepackage{amsmath,amssymb,amsfonts}
\usepackage{times}
\usepackage{graphicx}
\usepackage[caption=false,font=footnotesize]{subfig}
\usepackage{booktabs}
\usepackage{multirow}
\usepackage{array}
\usepackage{cite}
\usepackage{url}
\usepackage{balance}
\usepackage{siunitx}
\usepackage{enumitem}
\usepackage{tikz}
\usetikzlibrary{shapes.geometric,arrows.meta,fit,backgrounds}

% IEEE Publication Float Spacing Adjustments
\setlength{\textfloatsep}{2.5pt plus 1pt minus 1pt}
\setlength{\floatsep}{2.5pt plus 1pt minus 1pt}
\setlength{\intextsep}{2.5pt plus 1pt minus 1pt}
\setlength{\dbltextfloatsep}{3.5pt plus 1pt minus 1pt}
\setlength{\dblfloatsep}{2.5pt plus 1pt minus 1pt}
\setlength{\abovecaptionskip}{2pt}
\setlength{\belowcaptionskip}{2pt}

% Float placement parameters to avoid premature page ejection
\renewcommand{\topfraction}{0.92}
\renewcommand{\bottomfraction}{0.85}
\renewcommand{\textfraction}{0.08}
\renewcommand{\floatpagefraction}{0.85}
\renewcommand{\dbltopfraction}{0.92}
\renewcommand{\dblfloatpagefraction}{0.85}
\setcounter{topnumber}{4}
\setcounter{bottomnumber}{3}
\setcounter{totalnumber}{6}
\setcounter{dbltopnumber}{3}

% Compact list formatting to eliminate vertical gaps
\setlist{leftmargin=*,noitemsep,topsep=1.5pt,parsep=0pt,partopsep=0pt}
\def\BIBdecl{\setlength{\itemsep}{0.4pt plus 0.5pt minus 0.3pt}\parsep=0pt}

% Scientific Color Palette (IEEE Design System)
\definecolor{darknavy}{HTML}{202B3C}
\definecolor{acadblue}{HTML}{355C7D}
\definecolor{medblue}{HTML}{4C78A8}
\definecolor{mutedorange}{HTML}{E18D45}
\definecolor{mutedgreen}{HTML}{4A927A}
\definecolor{neutralgrey}{HTML}{737B87}
\definecolor{charcoal}{HTML}{343A40}
\definecolor{mutedwine}{HTML}{8B263E}

\begin{document}

\title{Mini-A-IPC: Design, Implementation, and Performance Evaluation of Classical and Adaptive Inter-Process Communication Mechanisms in Linux}

% Author Block: Student authors and course instructor metadata (no email field)
\author{
    \IEEEauthorblockN{Tanguturi Ram Sai (124CS0037)\IEEEauthorrefmark{1}, Ayush Naresh Gatla (124CS0038)\IEEEauthorrefmark{1}, and Dr.~N.~Srinivas Naik\IEEEauthorrefmark{2}}
    \IEEEauthorblockA{Department of Computer Science and Engineering\\
    Indian Institute of Information Technology Design and Manufacturing (IIITDM), Kurnool, India\\
    \IEEEauthorrefmark{1}Course: Operating Systems (CS303/CS353) \qquad \IEEEauthorrefmark{2}Course Instructor}
}

\maketitle

\begin{abstract}
Inter-Process Communication (IPC) is fundamental to concurrent multi-process systems and High-Performance Computing (HPC) software architectures. While the Linux operating system offers diverse IPC primitives---spanning byte streams, kernel message queues, and memory-mapped segments---selecting the optimal transport remains a non-trivial challenge dictated by message size, process topology, and virtualization. This paper presents Mini-A-IPC, an empirical evaluation platform and adaptive communication framework. We systematically evaluate six classical Linux IPC mechanisms (Anonymous Pipes, Named FIFOs, POSIX Shared Memory, System V Message Queues, POSIX Message Queues, and Unix Domain Sockets), implement a custom Linux character device mailbox (\texttt{/dev/aipc}) with an atomic reference-counting teardown barrier, and develop an adaptive transport framework incorporating static thresholding and dynamic EWMA telemetry with hysteresis. Across 3.43 million scheduled single-pair exchanges (2.73 million valid timed round-trip exchanges, yielding 5.46 million discrete directional message transmissions, and over 3.71 million total physically executed experimental cycles across all suites) on our primary Fedora 45 Linux environment (Kernel 7.3.3, KVM virtualization), we find that System V Message Queues achieve the lowest small-payload sample latency ($123.33\,\mu\text{s}$ at 64~B and $115.13\,\mu\text{s}$ at 256~B, descriptive lead under the KVM floor), intermediate transfers exhibit distinct empirical crossover optima (Unix Domain Sockets lead at 1~KB with $128.28\,\mu\text{s}$ and 64~KB with $123.53\,\mu\text{s}$, Named FIFOs lead at 4~KB with $119.44\,\mu\text{s}$, and Anonymous Pipes lead at 16~KB with $140.60\,\mu\text{s}$), and Unix Domain Sockets achieve superior scalability for large transfers ($449.64\,\mu\text{s}$ at 256~KB, $2.54\times$ faster than pipes, sustaining 1,304.44~MB/s wire bandwidth and 652.22~MB/s payload bandwidth; 1,244.01~MiB/s wire, paired $t$-test $p < 0.001$). In separate native driver validation, the custom kernel mailbox demonstrates an uncontended single-process loopback floor of $0.72\,\mu\text{s}$ ($P_{50} = 0.93\,\mu\text{s}$ at 1~KB; $2.95\text{--}3.35\,\mu\text{s}$ across-trial mean RTT for $\le 1$~KB transfers, scaling to $11.45\,\mu\text{s}$ at 16~KB), while multi-process evaluation on the primary KVM testbed was bounded by unprivileged credentials (\texttt{BLOCKED}) and the independent 95--130~$\mu\text{s}$ hypervisor scheduling floor. The static selector evaluates in 0.66~ns, achieving 57.14\% candidate accuracy between shared memory and pipes (14.29\% global accuracy across all transports; 27.68\% mean regret), while an offline-calibrated 3-tier policy provides an in-sample structural baseline.
\end{abstract}

\begin{IEEEkeywords}
Inter-process communication, Linux, operating systems, POSIX, System V, shared memory, message queues, adaptive transport selection, kernel modules, performance evaluation.
\end{IEEEkeywords}

\input{sections/01_introduction.tex}
\input{sections/02_related_work.tex}
\input{sections/03_architecture.tex}
\input{sections/04_ipc_implementation.tex}
\input{sections/05_adaptive_kernel.tex}
\input{sections/06_methodology.tex}
\input{sections/07_results.tex}
\input{sections/08_discussion.tex}
\input{sections/09_conclusion.tex}

\bibliographystyle{IEEEtran}
\bibliography{Mini-A-IPC-References}

\end{document}
